\section{Experiments and Results}
\label{sec:experiments}

We present our extensive multi-seed empirical evaluation designed to stress-test layer-wise model merging paradigms.

\subsection{Experimental Setup and Statistical Rigor}
We conduct all experiments using a pre-trained \textbf{CLIP ViT-B/32} backbone \cite{radford2021learning}. The network consists of 12 transformer blocks and 1 linear projection layer, resulting in $L = 13$ discrete parameter groups. 

\textbf{Task Experts:} We construct a multi-task benchmark consisting of 4 vision classification datasets:
\begin{enumerate}
    \item \textbf{MNIST} (Handwritten digits) \cite{lecun1998gradient}
    \item \textbf{FashionMNIST} (Fashion items) \cite{xiao2017fashion}
    \item \textbf{CIFAR-10} (Natural images) \cite{krizhevsky2009learning}
    \item \textbf{SVHN} (Street view house numbers) \cite{netzer2011reading}
\end{enumerate}

To ensure the highest standard of scientific and statistical rigor, we run our entire experimental pipeline across \textbf{3 independent random trials (using seeds 42, 100, and 2026)}. For each seed, we extract a unique disjoint subset of 512 images for expert training and another disjoint subset of 512 images for test-set evaluation. Experts are trained for 5 epochs using Adam with a learning rate of $10^{-5}$ on the CLIP backbone and $10^{-3}$ on task classification heads, ensuring they are properly converged and specialized. The mean expert test accuracies across the three seeds are: MNIST (\textbf{96.94\%}), FashionMNIST (\textbf{88.67\%}), CIFAR-10 (\textbf{88.93\%}), and SVHN (\textbf{85.81\%}), establishing a high-performing baseline for merging. Task vectors are extracted as $\tau^l_k = \theta^l_k - \theta^l_{\text{pre}}$.

\textbf{Coefficient Optimization:} For each trial, we construct an unlabeled calibration split of 256 images (64 images per task). We run both continuous optimization algorithms: the zero-order Adaptive 1+1 ES (for 500 steps) and the first-order Adam GD (for 200 steps) starting from uniform Task Arithmetic ($\lambda = 0.3$). Both optimizers successfully minimize the prediction entropy, reducing the calibration loss by over 30\% across trials.

\subsection{Comparison of Diagnostic Treatments}
Table \ref{tab:treatments} details the classification accuracies under all baselines and treatments, reporting the exact mean and standard deviation across the three independent random trials.

\begin{table*}[t]
\caption{Classification accuracies (Mean $\pm$ Std Dev \%) of model merging paradigms under diagnostic treatments across 3 independent random seeds on CLIP ViT-B/32. Task-wise spatial averaging on 1+1 ES (reducing parameter counts by 92.3\%) actually outperforms both the optimized layer-wise models and is statistically highly stable ($\pm 0.11\%$).}
\label{tab:treatments}
\centering
\resizebox{\textwidth}{!}{
\begin{tabular}{lcccccr}
\toprule
\textbf{Merging Paradigm / Treatment} & \textbf{MNIST} & \textbf{FashionMNIST} & \textbf{CIFAR10} & \textbf{SVHN} & \textbf{Average} & \textbf{Params / Task} \\
\midrule
\textbf{Task Arithmetic (Baseline, $\lambda=0.3$)} & $92.71 \pm 2.06$ & $81.64 \pm 1.00$ & $90.17 \pm 2.45$ & $73.24 \pm 1.54$ & $\mathbf{84.44 \pm 0.37}$ & 1 (Fixed) \\
\textbf{Optimized AdaMerging (1+1 ES)} & $93.62 \pm 2.02$ & $82.42 \pm 1.05$ & $89.13 \pm 2.40$ & $75.13 \pm 2.17$ & $\mathbf{85.07 \pm 0.47}$ & 13 (Optimized) \\
\textbf{Optimized AdaMerging (Adam GD)} & $92.84 \pm 1.53$ & $84.18 \pm 2.41$ & $89.84 \pm 1.25$ & $71.22 \pm 8.08$ & $\mathbf{84.52 \pm 1.57}$ & 13 (Optimized) \\
\midrule
\textbf{Intra-Task Layer Shuffling (1+1 ES)} & $92.58 \pm 2.72$ & $80.60 \pm 2.54$ & $84.77 \pm 1.38$ & $75.20 \pm 1.99$ & $\mathbf{83.28 \pm 1.26}$ & 13 (Shuffled) \\
\textbf{Intra-Task Layer Shuffling (Adam GD)} & $89.78 \pm 3.62$ & $81.71 \pm 4.72$ & $74.15 \pm 4.95$ & $70.70 \pm 11.03$ & $\mathbf{79.09 \pm 2.05}$ & 13 (Shuffled) \\
\textbf{Spatially Averaged (Spatial Mean - 1+1 ES)} & $93.62 \pm 1.78$ & $82.88 \pm 1.03$ & $86.65 \pm 1.79$ & $77.67 \pm 2.12$ & $\mathbf{85.21 \pm 0.11}$ & \textbf{1} (Mean scalar) \\
\textbf{Spatially Averaged (Spatial Mean - Adam GD)} & $92.38 \pm 1.52$ & $83.46 \pm 3.35$ & $79.49 \pm 4.62$ & $77.60 \pm 5.22$ & $\mathbf{83.24 \pm 0.99}$ & \textbf{1} (Mean scalar) \\
\bottomrule
\end{tabular}
}
\end{table*}

\begin{figure}[t]
\centering
\includegraphics[width=0.9\columnwidth]{results/fig1_treatments.png}
\caption{Classification accuracies under different treatments aggregated over 3 independent trials, showing standard error bars. Collapsing optimized coefficients into a flat task-wise Spatial Mean achieves competitive performance while eliminating overparameterization noise.}
\label{fig:treatments}
\end{figure}

The empirical results reveal a profound and unexpected dialectical interaction, which we term the \textbf{Overfitting-Optimizer Paradox}, where the apparent validity of "layer-specificity" is shown to depend fundamentally on the choice of optimization algorithm, transductive overfitting, and the underlying complexity of individual tasks:

\begin{itemize}
    \item \textbf{Zero-Order Optimization and the Layer-Specificity Illusion:} Under the derivative-free 1+1 ES optimizer, layer-specificity behaves entirely like a methodological illusion. Shuffling the optimized layer-wise coefficients (Intra-Task Layer Shuffling) results in only a minor, non-catastrophic performance decay from $85.07 \pm 0.47\%$ to $83.28 \pm 1.26\%$. More strikingly, collapsing the 13 parameters per task into a single task-wise scalar via Spatial Averaging (a 92.3\% reduction in parameters) actually improves average accuracy to $85.21 \pm 0.11\%$, while dramatically reducing variance across seeds. This empirical result suggests that under zero-order search, the optimized layer-specific coefficient variations do not reflect a physically meaningful, functional coordinate schedule. Instead, they are high-frequency optimization noise generated by random walk mutations in high-dimensional continuous space, which Spatial Averaging successfully smooths out, functioning as a powerful spatial regularizer that prevents parameter drift.
    
    \item \textbf{First-Order Optimization, Overfitting, and Delicate Layer-Specificity:} In sharp contrast, under standard first-order backpropagation-based Adam GD, layer-specific coordination appears to emerge as a highly functional mechanism. Shuffling optimized Adam coefficients collapses average performance from $84.52 \pm 1.57\%$ to $79.09 \pm 2.05\%$ (a significant $5.43\%$ drop), while replacing them with their Spatial Mean collapses average accuracy to $83.24 \pm 0.99\%$ (strictly worse than the unoptimized Task Arithmetic baseline of $84.44 \pm 0.37\%$). 
    
    However, a deeper methodological analysis reveals a profound paradox: \textbf{this apparent "functional reality" is actually an aggressive, delicate form of transductive overfitting} to the 256-image calibration set. First, Optimized Adam GD ($84.52 \pm 1.57\%$) fails to outperform the unoptimized Task Arithmetic baseline ($84.44 \pm 0.37\%$) on the unseen evaluation test set, while introducing massive instability and a 4x increase in standard deviation. Second, on CIFAR-10, the Optimized Adam GD model achieves $89.84 \pm 1.25\%$, which is strictly \textit{lower} than the unoptimized baseline ($90.17 \pm 2.45\%$). When optimizing 52 layer-wise parameters on a tiny calibration split using unconstrained first-order autograd, the optimizer successfully minimizes joint calibration entropy by finding a highly precise, delicate configuration of coefficients. This configuration is extremely sensitive to shuffling or spatial averaging, causing performance to collapse. But because this delicate structure is an artifact of overfitting to the transductive calibration statistics, it fails to translate to any generalizable performance improvements on unseen test samples. Thus, framing Adam GD's layer-specificity as a physically functional reality is overly generous; it is a manifestation of overparameterized transductive overfitting.
    
    \item \textbf{The SVHN Rescue vs. CIFAR-10 Collapse Trade-off and the Spatial Mean Illusion:} Looking beyond average performance reveals how multi-task metrics can mask catastrophic single-task behaviors, exposing the "Spatial Mean superiority" under 1+1 ES as a mathematical artifact of the SVHN rescue:
    \begin{itemize}
        \item On the highly complex \textbf{CIFAR-10} task, Spatial Averaging collapses performance by \textbf{10.35\%} under Adam GD (from $89.84\%$ to $79.49\%$) and by \textbf{2.48\%} under 1+1 ES (from $89.13\%$ to $86.65\%$). Under Adam GD, Intra-Task Shuffling collapses CIFAR-10 by a catastrophic \textbf{15.69\%} ($74.15\%$).
        \item On the sacrificial \textbf{SVHN} task, Spatial Averaging increases performance by \textbf{6.38\%} under Adam GD (from $71.22\%$ to $77.60\%$) and by \textbf{2.54\%} under 1+1 ES (from $75.13\%$ to $77.67\%$). 
    \end{itemize}
    This analysis proves that the Spatial Mean is not flatly "superior" or "sufficient." Rather, it acts as a coarse spatial regularizer. By collapsing the 13 layer coefficients to their mean, it successfully breaks the joint entropy minimization task-bias (which naturally sacrifices SVHN, as analyzed in Section 4.5) and stabilizes SVHN, but does so at the direct cost of destroying the delicate, functional layer-specific weight hierarchies of the more complex, non-linear CIFAR-10 task. This trade-off demonstrates that layer-specificity is physically functional and critical for complex tasks, and warns that simple average accuracies can easily mislead researchers in multi-task merging.
\end{itemize}

\subsection{Landscape Flatness Analysis}
To characterize the merging landscape, we inject relative Gaussian noise with scale factors $\gamma \in [0.05, 0.50]$ into the optimized parameters. Figure \ref{fig:noise} plots the average accuracy curves across seeds.

\begin{figure}[t]
\centering
\includegraphics[width=0.85\columnwidth]{results/fig2_noise_sensitivity.png}
\caption{Noise sensitivity analysis of optimized layer-wise merging coefficients under 1+1 ES and Adam GD (3 seeds). Shading represents standard deviation. Both optimization schedules tolerate up to 50\% relative noise with negligible performance decay, demonstrating an exceptionally flat optimization landscape.}
\label{fig:noise}
\end{figure}

Under 1+1 ES, adding 10\% relative noise ($\gamma = 0.10$) achieves **$85.19 \pm 0.34\%$** accuracy—nearly identical to the noise-free optimized model ($85.07\%$). At 50\% relative noise ($\gamma = 0.50$), performance remains highly robust at **$83.89 \pm 0.32\%$** (only a 1.18\% drop). Under Adam GD, injecting 50\% noise results in **$84.75 \pm 2.47\%$** accuracy. This extreme robustness to massive noise injection indicates that the exact values of learned layer coefficients are functionally irrelevant, confirming that layer-wise model merging is overparameterized and operates in an exceptionally flat loss basin.

\subsection{Representational Similarity (CKA) Results}
To track representational changes in activation space, we compute linear CKA at Layer 6 on CIFAR-10 inputs. Table \ref{tab:cka} outlines the CKA similarities between task experts and the merged models.

\begin{table*}[t]
\caption{Linear CKA similarity (Mean $\pm$ Std Dev \%) at Layer 6 on CIFAR-10 inputs across 3 seeds. While Spatially Averaged models exhibit marginally higher CKA values, the differences are statistically tiny (within standard deviation) and decoupling is observed as they do not translate to downstream classification accuracy.}
\label{tab:cka}
\centering
\resizebox{\textwidth}{!}{
\begin{tabular}{lcccc}
\toprule
\textbf{Expert (Ref)} & \textbf{Optimized (1+1 ES)} & \textbf{Averaged (1+1 ES)} & \textbf{Optimized (Adam GD)} & \textbf{Averaged (Adam GD)} \\
\midrule
\textbf{MNIST} & $0.9865 \pm 0.0039$ & $0.9896 \pm 0.0024$ & $0.9870 \pm 0.0021$ & $0.9891 \pm 0.0015$ \\
\textbf{FashionMNIST} & $0.9832 \pm 0.0096$ & $0.9915 \pm 0.0012$ & $0.9869 \pm 0.0076$ & $0.9896 \pm 0.0026$ \\
\textbf{CIFAR-10} & $0.9557 \pm 0.0307$ & $0.9633 \pm 0.0241$ & $0.9555 \pm 0.0302$ & $0.9598 \pm 0.0241$ \\
\textbf{SVHN} & $0.9799 \pm 0.0062$ & $0.9883 \pm 0.0022$ & $0.9830 \pm 0.0045$ & $0.9882 \pm 0.0008$ \\
\midrule
\textbf{Average CKA} & \textbf{0.9763} & \textbf{0.9832 (+0.0069)} & \textbf{0.9781} & \textbf{0.9817 (+0.0036)} \\
\bottomrule
\end{tabular}
}
\end{table*}

\begin{figure}[t]
\centering
\includegraphics[width=0.9\columnwidth]{results/fig3_cka.png}
\caption{Linear Centered Kernel Alignment (CKA) comparison across 3 seeds. Spatially averaged, flat task-wise merged models exhibit marginally higher CKA similarities on all tasks, though these differences are statistically small and do not guarantee downstream classification accuracy.}
\label{fig:cka}
\end{figure}

The CKA analysis reveals a highly consistent yet statistically marginal representational pattern: on average, \textbf{the spatially averaged model exhibits slightly higher activation similarity to original task experts (+0.0069 CKA for 1+1 ES and +0.0036 CKA for Adam GD) compared to the optimized model, though these individual differences are well within the measurement standard deviations and decouple from top-1 test accuracy under first-order gradient descent.} 

\textbf{The Optimization-Representation Trade-off and Spatial Regularization:} Under zero-order 1+1 ES optimization, this representational alignment corresponds to the superior generalization performance of the Spatial Mean. When optimizing $L \times K = 52$ continuous parameters on a calibration set of 256 images, the optimizers are highly susceptible to \textbf{transductive overfitting}. To minimize prediction entropy, the optimization process forces individual layer coefficients to overfit calibration samples, introducing spurious layer-by-layer variance that drifts intermediate activations away from the physical representations of original task experts. In contrast, Spatial Averaging collapses the parameters from 52 to just 4 task-wise scalars, acting as a powerful \textbf{spatial regularizer}. By enforcing spatial uniformity across layers, Spatial Averaging prevents parameter drift, mitigates transductive overfitting, and preserves activation-level representations, leading to both superior representational alignment and superior test generalization.

\textbf{The CKA-Accuracy Discrepancy and Representation Limits:}
However, looking closely at first-order Adam GD optimization reveals a profound discrepancy: while the Spatially Averaged model has a slightly higher average activation similarity to original experts, this representational alignment directly contradicts downstream classification accuracy. For instance, on CIFAR-10 under Adam GD, the Spatially Averaged model has a higher CKA ($0.9598 \pm 0.0241$) than the Optimized model ($0.9555 \pm 0.0302$), yet its classification accuracy collapsed catastrophically by $10.35\%$ (from $89.84\%$ to $79.49\%$). 

This exposes a profound methodological limitation of activation-level similarity metrics like CKA: \textbf{high-level linear kernel alignment is a poor predictor of downstream classification performance in the fine-grained model merging regime}. Because neural networks are highly non-linear, slight coordinate adjustments or decision boundary shifts in weight-space can cause catastrophic misclassifications of target samples even while keeping the overall activation subspace highly correlated ($>0.95$ CKA). While spatial regularization successfully preserves a coarse, global activation alignment, first-order gradient descent constructs highly specialized weight-space coordinations that are critical to maintaining decision boundary integrity. Researchers must be highly cautious when using representational similarity as a proxy for physical task performance.

\subsection{Methodological Discussions: Optimizer Confounding, Task Complexity, and Joint Entropy Bias}

We provide a self-critical, rigorous analysis of three key methodological factors:

\textbf{1. Optimizer Confounding and the Overfitting Paradox:}
By implementing first-order backpropagation-based Adam GD alongside zero-order 1+1 ES, we isolate the critical interaction between optimizer capabilities, parameterization, and transductive overfitting. Under zero-order 1+1 ES search, the overparameterized coefficient schedule (52 parameters) tuned on a small calibration set of 256 images is highly susceptible to transductive overfitting, which takes the form of high-frequency optimization noise easily resolved and smoothed out by Spatial Averaging (acting as a coarse regularizer). 

However, when exact first-order gradients are available (Adam GD), the optimizer successfully discovers highly precise, delicate layer-specific weight balances that minimize joint calibration entropy. Shuffling these coefficients breaks these delicate structures, causing a 5.43\% drop in average accuracy and a catastrophic 15.69\% collapse on CIFAR-10. This creates a powerful illusion of highly functional layer-specific coordination. Yet, because this configuration is an artifact of transductive overfitting, it fails to outperform the unoptimized Task Arithmetic baseline on unseen test samples ($84.52\%$ vs. $84.44\%$) while multiplying the standard deviation across seeds by 4x ($\pm 1.57\%$ vs. $\pm 0.37\%$). This demonstrates that layer-specificity under unconstrained first-order gradient descent is a manifestation of delicate, overfit transductive drift on the calibration set. Furthermore, to investigate whether standard optimizer tuning can mitigate this transductive overfitting without explicit regularization, we conduct a systematic sweep of the learning rate ($\eta$) under unconstrained Adam GD in Appendix \ref{sec:appendix_learning_rate}. We find that although lower learning rates restrict parameter drift and act as a crude form of early-stopping regularization (preserving CIFAR-10 performance), they also limit convergence, validating our proposed explicit coefficient regularization as a superior and more targeted solution.

\textbf{2. Task Complexity and Layer Coordination:}
While average performance remains highly functional under shuffling, a closer look at the hardest task (SVHN, expert accuracy ~85.81\%) reveals a critical exception. Shuffling Adam GD's coefficients causes SVHN performance to drop to **$70.70 \pm 11.03\%$**, exhibiting massive variance. This indicates that \textbf{layer-wise coordination does possess physical, functional importance for harder, more complex, or highly non-linear tasks} that reside further from the pre-trained initialization. For simpler tasks (MNIST/FashionMNIST), the task vectors lie extremely close to the pre-trained CLIP base, rendering layer-wise coefficients redundant. Thus, layer-specificity is an illusion for simple domains, but becomes functional in harder ones.

\textbf{3. Joint Entropy Minimization Task-Bias:}
Our empirical results expose a critical, unformulated multi-task bias in joint entropy minimization objectives. Under Adam GD, SVHN test accuracy collapses to **$71.22 \pm 8.08\%$**—lower than both Task Arithmetic (73.24\%) and 1+1 ES (75.13%). Because SVHN is highly complex, its prediction entropy on the calibration set is naturally much higher than MNIST's. To minimize the joint sum, the optimizer trades off SVHN's parameters to eke out marginal entropy reductions on simpler, low-entropy tasks. This bias sacrifices the harder tasks to minimize joint loss, degrading multi-task capability. We recommend that future test-time adaptation objectives adopt weighted or temperature-scaled entropy formulations to prevent this sacrificial task behavior. To evaluate this recommendation, we present a pilot study in Appendix E formulating and executing a Scale-Normalized Weighted Joint Entropy objective, weighting each task's entropy by the inverse of its baseline uniform task arithmetic entropy. Our pilot results confirm that scale-normalization successfully resolves this task-bias and improves test performance (retaining 85.84\% average accuracy).

\textbf{4. Proposed Solution: Explicit Coefficient Regularization:}
To mitigate transductive overfitting, reduce optimization instability (the 4x increase in seed-to-seed variance), and prevent performance collapse on complex tasks (such as CIFAR-10), we propose that future test-time adaptive model merging frameworks incorporate explicit coefficient regularization into their objectives. Rather than running unconstrained test-time adaptation, the calibration loss should be augmented with an $L_2$ regularization penalty or proximity penalty:
\begin{equation}
    \mathcal{L}_{\text{reg}} = \mathcal{L}_{\text{calibration}} + \beta \sum_{k=1}^K \sum_{l=1}^L ||\lambda^l_k - \lambda_{\text{init}}||^2_2
\end{equation}
where $\lambda_{\text{init}}$ is the uniform initialization (e.g., 0.3 for Task Arithmetic) and $\beta > 0$ is a regularization hyperparameter. This penalty restricts the coefficients from drifting excessively from their stable, spatially uniform baseline, preventing the optimizer from finding delicate, overfit parameter configurations on small calibration samples. By enforcing proximity to the flat baseline, coefficient regularization can stabilize the optimization, reduce seed-to-seed variance, and preserve the delicate representational hierarchies of complex tasks. 

To verify this, we conduct an empirical pilot sweep of $\beta \in [0.0, 1.0]$ in Appendix B, proving that proximity regularization stabilizes test performance at a peak average accuracy of 86.57\% ($\beta = 0.5$). Furthermore, in Appendix F, we compare Proximity Regularization with standard optimizer-level weight decay (which pulls coefficients to 0.0), demonstrating that standard weight decay causes a catastrophic collapse of SVHN performance (to 52.54\%) whereas Proximity Regularization successfully preserves and stabilizes all experts.
